\documentclass[a4paper]{article}
\usepackage{amsfonts}
\usepackage{amsmath}
\usepackage{amssymb}
\usepackage{graphicx}     
\usepackage{cancel}

% Command definities van frequent gebruikte commandos
\newcommand{\N}{\mathbb{N}}
\newcommand{\Q}{\mathbb{Q}}
\newcommand{\R}{\mathbb{R}}
\newcommand{\C}{\mathbb{C}}
\newcommand{\Bgsin}{\operatorname{Bgsin}}
\newcommand{\Bgcos}{\operatorname{Bgcos}}
\newcommand{\Bgtan}{\operatorname{Bgtan}}
\newcommand{\cis}{\operatorname{cis}}
\newcommand{\Limit}{\lim\limits_}
\newcommand{\Int}{\displaystyle\int}

\begin{document}
  
\noindent \large Auteur: Wout Vanderborght\\
\noindent \large Studierichting: Informatica\\
\noindent \large Oefeningenreeks 4A nr. 34.7\\

\medskip

\normalsize

$ \Int \dfrac{e^{3x}}{e^{3x}+ 1}dx $ \\

Stel $t = e^{3x}+1 \Rightarrow dt = 3e^{3x}dx \Leftrightarrow \dfrac{dt}{3} = e^{3x}dx $\\

$ = \dfrac{1}{3} \Int \dfrac{dt}{t} $ \\

$ = \dfrac{1}{3} \ln|t| + C $ \\

$ = \dfrac{1}{3} \ln|e^{3x}+1| + C $ \\

\begin{thebibliography}{999}
%\bibitem{a} Referentie
\end{thebibliography}
\end{document} 
